\section{How Far a Null Climbs}
\label{sec:ch14-how-far-a-null-climbs}

\index{null propagation!the rule, restated across a seam|(}%
Chapter~\ref{ch:schema-design} quoted the rule and walked it inside one
service, and it has not changed~\autocite{graphqlspec2025}:

\begin{quote}
Since Non-Null response positions cannot be null, execution errors are
propagated to be handled by the parent response position. If the parent
response position may be null then it resolves to null, otherwise if it is a
Non-Null type, the execution error is further propagated to its parent
response position.
\end{quote}

\index{specification!where the climb stops being survivable}%
What chapter~\ref{ch:schema-design} did not need is the sentence that closes
that section~\autocite{graphqlspec2025}:

\begin{quote}
If every response position from the root of the request to the source of the
execution error has a Non-Null type, then the \code{data} entry in the
execution result should be null.
\end{quote}

\index{null propagation!nothing federated about it}%
That is the whole mechanism behind section~\ref{sec:ch14-the-response-that-is-not-there}.
There is nothing federated in either paragraph. A router that merges four
services into one response has not changed the rule; it has changed how many
positions the walk passes through and who wrote each of them.

\subsection{The Walk, Three Times}

\index{null propagation!where each of the three stops}%
Take the three requests this chapter has sent and walk each one upward from
where it failed. Figure~\ref{fig:ch14-the-climb-crosses} draws all three
against the same tree.

\begin{figure}[htbp]
  \centering
  % The same failure, raised at three positions in one supergraph, and the three
% different amounts of the response it destroys.
%
% Same idiom as figures/tikz/ch01-one-field.tex, which fixes it: 0.4pt strokes,
% the Straight Barb arrowhead, sans-serif scriptsize labels, square corners, no
% fills, black on white. Closest relative in the book is
% figures/tikz/ch04-null-climbs.tex, which draws one climb inside one service;
% this is the federated version of that picture and keeps that file's
% conventions on purpose: a box per response position, its GraphQL type inside
% it where the type is what decides the outcome, solid for a position present
% in the response and dashed for one that is not there at all.
%
% Three columns rather than three figures, because the failure is identical in
% all three and only the wrapper differs, which is only visible side by side.
%
% An earlier draft drew the seam as one horizontal rule across all three
% columns. That is wrong and it is worth recording why, because the rule looked
% right: in the third column the stopped service answers a ROOT field, one step
% under data, while a healthy service answers its sibling at the same depth. A
% horizontal line cannot separate two siblings, so the seam is carried by the
% caption and by the leaf labels instead.
%
% Bare tikzpicture (house style): the figure environment, caption and label
% live at the call site in the section file.
\begin{tikzpicture}[
  x=1mm, y=1mm,
  every node/.append style={font=\sffamily\scriptsize},
  posbox/.style={draw, line width=0.4pt, minimum width=24mm,
    minimum height=8mm, align=center, inner sep=1pt},
  smallbox/.style={posbox, minimum width=21mm},
  lostbox/.style={posbox, dashed},
  lostsmall/.style={smallbox, dashed},
  link/.style={draw, line width=0.4pt},
  lostlink/.style={draw, line width=0.4pt, dashed},
  climb/.style={draw, line width=0.4pt,
    -{Straight Barb[length=1.4mm, width=1.6mm]}},
  lane/.style={anchor=north, align=center, font=\sffamily\scriptsize\itshape,
    text width=42mm},
  note/.style={anchor=north, align=center, inner sep=2pt, text width=42mm},
  leaf/.style={anchor=north, align=center, inner sep=1pt, text width=42mm,
    font=\sffamily\scriptsize\itshape},
  % The two tags in column 3 sit under boxes half a column apart, so they take
  % the width of the box rather than of the column.
  leaftag/.style={leaf, text width=21mm},
]

% =================================================================== column 1
\node[posbox] (a-root)  at (18,64) {data};
\node[posbox] (a-sess)  at (18,48) {sessions};
\node[posbox] (a-nodes) at (18,32) {nodes\\{}[Session!]};
\node[posbox] (a-item)  at (18,16) {nodes[0]\\Session!};
\node[posbox] (a-field) at (18,0)  {averageScore\\Float};

\draw[link] (a-root.south)  -- (a-sess.north);
\draw[link] (a-sess.south)  -- (a-nodes.north);
\draw[link] (a-nodes.south) -- (a-item.north);
\draw[link] (a-item.south)  -- (a-field.north);

\node[leaf] at (18,-6) {the stopped service answers here};
\node[lane] at (18,-13) {one field};
\node[note] at (18,-19)
  {the position may be null,\\so the walk never starts};

% =================================================================== column 2
\node[posbox]  (b-root)  at (68,64) {data};
\node[posbox]  (b-sess)  at (68,48) {sessions};
\node[posbox]  (b-nodes) at (68,32) {nodes\\{}[Session!]};
\node[lostbox] (b-item)  at (68,16) {nodes[0]\\Session!};
\node[lostbox] (b-field) at (68,0)  {ratingCount\\Int!};

\draw[link] (b-root.south) -- (b-sess.north);
\draw[link] (b-sess.south) -- (b-nodes.north);

\draw[climb] (b-field.north) -- (b-item.south);
\draw[climb] (b-item.north)  -- (b-nodes.south);

\node[leaf] at (68,-6) {the stopped service answers here};
\node[lane] at (68,-13) {four sessions};
\node[note] at (68,-19)
  {two positions cannot hold it,\\and the list that can is the\\whole page};

% =================================================================== column 3
% The root field and its sibling are at the same depth, which is the whole
% reason this column loses more than the other two.
% Solid, like the nodes box in column 2 and for the same reason: the position
% is in the response holding a null. Only what it would have contained is gone.
\node[posbox]   (c-root)  at (120,64) {data};
\node[lostsmall](c-field) at (107,46) {ratingCount\\Int!};
\node[lostsmall](c-sess)  at (133,46) {sessions};

\draw[climb]    (c-field.north) -- (107,56) -- (c-root.south);
\draw[lostlink] (c-root.south)  -- (133,56) -- (c-sess.north);

\node[leaftag] at (107,40) {stopped};
\node[leaftag] at (133,40) {answered};
\node[lane] at (120,-13) {the whole response};
\node[note] at (120,-19)
  {a root field's parent is the\\response, and the sibling\\beside it goes too};

% ----------------------------------------------------------------- legend
% Below the deepest note rather than beside it: the three notes run to three
% lines and a legend level with them collides with all of them.
\node[posbox, minimum width=6mm, minimum height=3.5mm] at (4,-38) {};
\node[note, anchor=west, text width=44mm] at (9,-35.5)
  {present in the response};

\node[lostbox, minimum width=6mm, minimum height=3.5mm] at (62,-38) {};
\node[note, anchor=west, text width=44mm] at (67,-35.5)
  {not in the response at all};

\end{tikzpicture}

  \caption{One stopped service, three requests, and the three different amounts
  of the response it destroys. What decides each one is not the failure, which
  is identical in all three, but the distance up to the first position whose
  type is allowed to hold a null. On the left that is no distance at all. In
  the middle it is two positions, and the list that finally absorbs it is the
  whole page. On the right the failing field is a root field, so the only
  position above it is the response, and the \code{sessions} beside it is a
  sibling at the same depth that a running service answered correctly and that
  goes anyway. A solid box is a position the response still has, whether or not
  it holds anything: \code{nodes} in the middle and \code{data} on the right
  are both present and both null.}
  \label{fig:ch14-the-climb-crosses}
\end{figure}

\index{null propagation!the nullable field}%
\code{averageScore} is \code{Float}. The router cannot fill it, the position may
be null, and the walk ends before it starts. Four sessions, four nulls.

\index{null propagation!through a list of non-null items}%
\code{ratingCount} is \code{Int!} and cannot hold it, so the error goes to its
parent, the \code{Session} at index zero. \code{nodes} is typed
\code{[Session!]}, so the items may not be null either and the error keeps
going. The parent of an item is the list, the list is nullable, and it stops
there. The three sessions that had nothing wrong with them are gone because
they were in the same list as the one that did.

\index{null propagation!from a root field}%
\code{ratingCount} at the root is \code{Int!} and its parent is the response.
There is no position between them.

\subsection{Two Things the Walk Explains}

\index{errors!fewer errors for more damage}%
The first is the error count, which runs the wrong way round. With Speakers
stopped, a request for \code{speaker \{ name \}} across four sessions comes back
with five errors: the fetch failure, and one
\enquote{Cannot return null for non-nullable field} per session, at paths zero
through three. The request that lost all four sessions to
\code{ratingCount} carries two. Once the list is gone there are no positions
left to report, so the response that lost the most says the least. Sorting an
error log by count puts these exactly the wrong way up.

\index{null propagation!the sibling branch is untouched}%
The second is what survives. Ask for \code{pageInfo} beside the broken
\code{nodes}, and the \code{data} half of the answer, under the same two errors
as before, is this:

\begin{minted}[fontsize=\footnotesize,breakanywhere]{json}
{"data":{"sessions":{"nodes":null,"pageInfo":{"hasNextPage":false,"endCursor":"Mw=="}}}}
\end{minted}

\index{null propagation!stops at the common parent}%
The cursor is intact. \code{pageInfo} is a sibling of \code{nodes} under
\code{sessions}, and the climb stopped at \code{nodes} without ever reaching
their common parent. So a client can be handed a page whose contents are null
and whose pagination is fine, which is a shape worth knowing about before a
client tries to advance through it.

\index{null propagation!edges bubbles one level higher}%
And asking through \code{edges} rather than \code{nodes} loses the same four
sessions one position higher up, because \code{SessionsEdge.node} is
\code{Session!} and absorbs nothing on the way past. Two spellings of the same
query, two different things null in the response.
\index{null propagation!the rule, restated across a seam|)}%
